\section*{Capítulo \ref{ch:b1:stereochemistry-in-depth} --- Estereoquímica aprofundada}

\begin{solution}{exo:b1:stereochemistry-in-depth:1}
\omterm{def:b1:stereochemistry-in-depth:isomers}{Isômeros constitucionais}; \omterm{def:b1:stereochemistry-in-depth:enantiomers}{enantiômeros}; \omterm{def:b1:stereochemistry-in-depth:enantiomers}{diastereoisômeros}
(\omterm{def:b1:stereochemistry-in-depth:isomers}{estereoisômeros} que não são imagens especulares); duas \omterm{def:b1:stereochemistry-in-depth:isomers}{conformações} de
um mesmo composto; \omterm{def:b1:stereochemistry-in-depth:enantiomers}{diastereoisômeros} (\cip{R,S} é a forma meso).
\end{solution}

\begin{solution}{exo:b1:stereochemistry-in-depth:2}
\ce{-OH} $>$ \ce{-NH2} $>$ \ce{-CH3} $>$ \ce{-H} (O, N, C, H).
\ce{-COOH} (O,O,O) $>$ \ce{-CHO} (O,O,H) $>$ \ce{-CH2OH} (O,H,H) $>$
\ce{-CH(CH3)2} (C,C,H). \ce{-Br} $>$ \ce{-Cl} $>$ \ce{-CCl3} (Cl,Cl,Cl)
$>$ \ce{-CH2Cl} (Cl,H,H): o primeiro átomo decide antes de seus vizinhos.
\end{solution}

\begin{solution}{exo:b1:stereochemistry-in-depth:3}
2,3-Diclorobutano: dois centros equivalentes, \cip{R,R}, \cip{S,S} e o
meso \cip{R,S}: três. 2-Bromo-3-clorobutano: dois centros diferentes,
quatro \omterm{def:b1:stereochemistry-in-depth:isomers}{estereoisômeros}, nenhuma forma meso. Pentano-2,4-diol: três, sendo
\cip{2R,4S} o meso.
\end{solution}

\begin{solution}{exo:b1:stereochemistry-in-depth:4}
Em C5: isopropenila (C,C,C, a ligação dupla contando duas vezes) $>$ C6
$>$ C4 $>$ H. C6 e C4 levam ambos (C,H,H); um passo adiante, C6 leva ao
carbono da carbonila (O,O,C) e C4 a C3 (C,C,H): C6 vence. Com o grupo
isopropenila na frente e H atrás, a ordem isopropenila $\to$ C6 $\to$ C4
gira no sentido horário, como desenhado: \cip{R}, a carvona da hortelã; a
de cunha tracejada é
\cip{S}.
\end{solution}

\begin{solution}{exo:b1:stereochemistry-in-depth:5}
$\theta = 0$: metilas eclipsados, \qty{16.5}{kJ/mol}; $60^\circ$: gauche,
cerca de \qty{2.9}{kJ/mol} (mínimo de 2.8 em $62^\circ$); $120^\circ$:
metila eclipsando H, \qty{15.2}{kJ/mol}; $180^\circ$: anti, 0. Com
energias iguais, duas \omterm{def:b1:stereochemistry-in-depth:conformations}{conformações gauche} para uma anti: dois terços
gauche.
\end{solution}

\begin{solution}{exo:b1:stereochemistry-in-depth:6}
Cadeia vertical com \ce{COOH} em cima e \ce{CH3} embaixo; prioridades
$\ce{NH2} > \ce{COOH} > \ce{CH3} > \ce{H}$. Pondo \ce{NH2} à esquerda e H
à direita, $\ce{NH2} \to \ce{COOH} \to \ce{CH3}$ gira no sentido horário
como visto; H está na horizontal (para o leitor), de modo que o
descritor se inverte: \cip{S}. O grupo amino fica à esquerda.
\end{solution}

\begin{solution}{exo:b1:stereochemistry-in-depth:7}
\textit{cis}: em carbonos adjacentes, uma ligação para cima e outra para
baixo entre as posições axiais e \omterm{def:b1:stereochemistry-in-depth:conformations}{equatoriais} dá axial--equatorial nas
duas \omterm{def:b1:stereochemistry-in-depth:conformations}{cadeiras}. \textit{trans}: diequatorial ou diaxial. Na \omterm{def:b1:stereochemistry-in-depth:conformations}{cadeira}
diequatorial, os dois metilas estão em gauche um com o outro (uma
interação); o isômero \textit{cis} tem essa interação mais um metila
\omterm{def:b1:stereochemistry-in-depth:conformations}{axial}, cerca de \qty{5.7}{kJ/mol} a mais. O isômero \textit{trans} é o
mais estável.
\end{solution}

\begin{solution}{exo:b1:stereochemistry-in-depth:8}
$[\alpha] = 13.2/(2.00 \times 0.100) = +66.0^\circ$. Um tubo de
\qty{1.00}{dm}: $+6.60^\circ$. O \omterm{def:b1:stereochemistry-in-depth:enantiomers}{enantiômero}: $-13.2^\circ$ no tubo de
\qty{2.00}{dm}.
\end{solution}

\begin{solution}{exo:b1:stereochemistry-in-depth:9}
\qty{25}{\celsius}: $\exp(5700/(8.314 \times 298.15)) = 10.0$, fração
$10.0/11.0 = \qty{91}{\percent}$. \qty{-50}{\celsius}: $\exp(5700/(8.314
\times 223.15)) = 21.6$, \qty{96}{\percent}.
\end{solution}

\begin{solution}{exo:b1:stereochemistry-in-depth:10}
No etano, os seis hidrogênios são iguais: toda \omterm{def:b1:stereochemistry-in-depth:conformations}{conformação alternada} é a
mesma. No butano, o metila de trás pode ficar oposto ao da frente (anti)
ou ao lado dele ($\pm 60^\circ$, gauche), onde os dois metilas se
comprimem. Populações: anti 1, gauche $2\exp(-2.8/2.48) = 0.65$: cerca
de \qty{61}{\percent} anti e \qty{39}{\percent} gauche.
\end{solution}

\begin{solution}{exo:b1:stereochemistry-in-depth:11}
C2 e C4 são estereogênicos. C3 leva H, OH e dois ramos de constituição
idêntica; eles só diferem quando C2 e C4 têm descritores opostos, e só
então C3 é estereogênico (pseudoassimétrico). Isômeros:
\cip{2R,4R} e \cip{2S,4S}, um par de \omterm{def:b1:stereochemistry-in-depth:enantiomers}{enantiômeros} (C3 não estereogênico);
\cip{2R,4S} com os dois arranjos em C3, duas formas meso. Quatro ao todo.
\end{solution}

\begin{solution}{exo:b1:stereochemistry-in-depth:12}
$[\alpha]_{\text{amostra}} = 2.31/(1.00 \times 0.050) = +46.2^\circ$, e
$2x - 1 = 46.2/66.0 = 0.700$: $x = 0.850$, \qty{85}{\percent} do
\omterm{def:b1:stereochemistry-in-depth:enantiomers}{enantiômero} $(+)$ e \qty{15}{\percent} do $(-)$.
\end{solution}

\begin{solution}{pb:b1:stereochemistry-in-depth:1}
\textbf{1.} Anel de ciclo-hexano: C1 leva o OH, C2 o grupo isopropila,
C5 o metila; C1, C2 e C5 são estereogênicos.
\textbf{2.} $2^3 = 8$, quatro pares de \omterm{def:b1:stereochemistry-in-depth:enantiomers}{enantiômeros}.
\textbf{3.} Nenhum arranjo tem plano especular: os três substituintes
são todos diferentes.
\textbf{4.} \ce{OH} $>$ C2 (C,C,H) $>$ C6 (C,H,H) $>$ H.
\textbf{5.} C1 (O,C,H) $>$ isopropila (C,C,H) $>$ C3 (C,H,H) $>$ H.
\textbf{6.} C6 e C4 ambos (C,H,H); em seguida, C6 leva a C1 (O,C,H), C4 a
C3 (C,H,H): C6 $>$ C4 $>$ \ce{CH3} $>$ H.
\textbf{7.} \cip{1S,2R,5S}.
\textbf{8.} Em uma \omterm{def:b1:stereochemistry-in-depth:conformations}{cadeira}, duas ligações \omterm{def:b1:stereochemistry-in-depth:conformations}{equatoriais} em carbonos
adjacentes (C1, C2) apontam uma para cima e outra para baixo:
\textit{trans}; nos carbonos 1 e 3 do anel (C1 e C5, passando por C6),
as duas ligações \omterm{def:b1:stereochemistry-in-depth:conformations}{equatoriais} apontam no mesmo sentido: \textit{cis}. Esse
é o arranjo do mentol.
\textbf{9.} Os três ficam \omterm{def:b1:stereochemistry-in-depth:conformations}{axiais}.
\textbf{10.} Na \omterm{def:b1:stereochemistry-in-depth:conformations}{cadeira} invertida, só o metila custaria cerca de
\qty{5.7}{kJ/mol}, e o grupo isopropila, maior, mais; o OH acrescenta um
custo próprio menor, e o OH e o metila \omterm{def:b1:stereochemistry-in-depth:conformations}{axiais}, nos carbonos 1 e 3 e na
mesma face, também se comprimiriam. Isso dá bem mais de \qty{12}{kJ/mol},
uma razão de mais de $\exp(12000/2479) \approx 130$ para um: o mentol
está quase inteiramente na \omterm{def:b1:stereochemistry-in-depth:conformations}{cadeira} toda \omterm{def:b1:stereochemistry-in-depth:conformations}{equatorial}.
\textbf{11.} Um \omterm{def:b1:stereochemistry-in-depth:enantiomers}{diastereoisômero}: só um dos três centros difere.
\textbf{12.} Com o isopropila e o metila \omterm{def:b1:stereochemistry-in-depth:conformations}{equatoriais}, o OH fica \omterm{def:b1:stereochemistry-in-depth:conformations}{axial}.
\textbf{13.} Os \omterm{def:b1:stereochemistry-in-depth:enantiomers}{diastereoisômeros} diferem em todas as suas propriedades
(aqui, as \omterm{def:b1:intermolecular-forces-solvents:hbond}{ligações de hidrogênio} de um OH \omterm{def:b1:stereochemistry-in-depth:conformations}{axial} ou \omterm{def:b1:stereochemistry-in-depth:conformations}{equatorial}); os
\omterm{def:b1:stereochemistry-in-depth:enantiomers}{enantiômeros} só diferem diante de parceiros \omterm{def:b1:stereochemistry-in-depth:stereocentre}{quirais}.
\textbf{14.} $[\alpha] = -2.46/(1.00 \times 0.0500) = -49.2^\circ$, dentro
do intervalo do manual.
% ledger: alpha:l-menthol-range
\textbf{15.} $+49.2^\circ$; $0$.
\textbf{16.} $-1.97/0.0500 = -39.4^\circ$.
\textbf{17.} $\alpha = \ell c[\alpha]_{(-)}\bigl(x - (1 - x)\bigr) =
\ell c[\alpha]_{(-)}(2x - 1)$.
\textbf{18.} $x = \frac12(1 + \alpha/\alpha_{\text{ref}})$ com $\ell$ e $c$ iguais.
\textbf{19.} Sim: qualquer outro composto opticamente ativo acrescenta
seu próprio termo (\omterm{prop:b1:stereochemistry-in-depth:biot}{lei de Biot}). É preciso verificar a pureza química por
outro método (cromatografia) antes de ler a rotação como uma composição.
\textbf{20.} $x = \frac12(1 + 1.97/2.46) = 0.900$.
\textbf{21.} $x = \frac12(1 + 2.40/2.46) = 0.988$.
\textbf{22.} A segunda (\qty{98.8}{\percent}).
\textbf{23.} $-3.94^\circ$; o mesmo erro de leitura passa a ser uma
fração menor do ângulo.
\textbf{24.} \textbf{\qty{90.0}{\percent} de $(-)$-mentol} (e
\qty{10.0}{\percent} de seu \omterm{def:b1:stereochemistry-in-depth:enantiomers}{enantiômero}).
\end{solution}
